\section{Working physical hypothesis}
\label{sec:hypothesis}

The preceding theory suggests a hypothesis that is stronger than a statement of
numerical cancellation and weaker than a claimed universal theorem.

\begin{keypoint}[Channel-selective protection hypothesis]
For a low-energy electron--phonon process, the leading error of static DFPT is
controlled by (i) the decomposition of the self-consistent phonon perturbation
into charge, internal, nonlocal, and interaction-modulation operators; (ii) the
many-body amplification in each operator channel; and (iii) the momentum and
frequency path of the response.  Strong correlation changes the error mainly when
the phonon has appreciable weight in an unprotected channel with a large or
dynamic susceptibility.
\end{keypoint}

Write the projected perturbation in an orthonormal operator basis,
\begin{equation}
  D_{\nu\vq}^{\ks}=\sum_a d_{a,\nu\vq}O_a,
  \qquad
  \Tr(O_a^\dagger O_b)=\delta_{ab}.
  \label{eq:operatorbasis}
\end{equation}
The corresponding many-body perturbation is represented by a channel kernel
\begin{equation}
  D_{\nu\vq}^{\mb}(\omega)
  =\sum_{ab}\mathcal K_{ab}(\vq,\omega)
  d_{b,\nu\vq}O_a.
  \label{eq:channelkernel}
\end{equation}
The first implementation may retain only diagonal or small-block kernels, but the
general expression allows channel mixing.  The hypothesis predicts the following.

\begin{enumerate}
  \item The strict uniform common-charge direction has an exact consistency
  constraint: it is an energy-zero shift and cannot acquire an arbitrary
  quasiparticle renormalization.
  \item The finite-$q$ density channel is not exactly protected, but UEG evidence
  suggests that it may often be represented by a compact static kernel.
  \item Traceless on-site orbital operators, hopping/hybridization operators, and
  interaction derivatives are not constrained by charge conservation and may
  receive large, mode-specific corrections.
  \item Static and dynamic validity are separate.  A channel may have
  $\mathcal K_a(\vq,0)\simeq1$ while varying rapidly across
  $\omega_{\nu\vq}$.
\end{enumerate}

\subsection{Measurable channel weights}

A simple Frobenius norm of an operator can overemphasize matrix elements that do
not connect the relevant electronic states.  For a target energy window
$\mathcal W$ and nonnegative sampling weights $w_{n\vk}$, define a
Fermi-window channel strength
\begin{equation}
  S_{a,\nu\vq}
  =\sum_{mn\vk\in\mathcal W}
   w_{n\vk}w_{m,\vk+\vq}
   \left|
   \langle\psi_{m,\vk+\vq}|
   d_{a,\nu\vq}O_a
   |\psi_{n\vk}\rangle
   \right|^2.
  \label{eq:channelstrength}
\end{equation}
Normalize $w_a=S_a/\sum_bS_b$ for interpretation.  These weights answer which
operators the actual low-energy scattering process samples, not merely which
matrix block is largest in a localized representation.

A correction-risk score should combine channel weight and amplification.  A
minimal form is
\begin{equation}
  R_{\nu\vq}
  =\left[
   \sum_a w_{a,\nu\vq}
   \left|\mathcal K_a(\vq,0)-1\right|^2
  \right]^{1/2}
  +\lambda_{\mathrm{dyn}}\max_a\eta^{\mathrm{dyn}}_{a,\nu\vq},
  \label{eq:risk}
\end{equation}
where
\begin{equation}
  \eta^{\mathrm{dyn}}_{a,\nu\vq}
  =\omega_{\nu\vq}
   \left|
   \partial_\omega\ln\mathcal K_a(\vq,\omega)
   \right|_{\omega=0}.
  \label{eq:dynamicrisk}
\end{equation}
The threshold and $\lambda_{\mathrm{dyn}}$ are calibration parameters, not
universal constants.

\subsection{Falsification criteria}

The hypothesis must be rejected or revised if any of the following occur across a
held-out benchmark set:

\begin{itemize}
  \item channel weights do not correlate with the size or sign of the
  beyond-DFPT correction;
  \item nominally charge-dominated modes show large, uncontrolled corrections
  that cannot be explained by finite-$q$, local-field, or dynamic effects;
  \item the kernel requires essentially independent parameters for every
  band--momentum matrix element, eliminating the proposed compression; or
  \item dynamic effects dominate so broadly that a static risk gate cannot
  identify a useful low-cost regime.
\end{itemize}

The theory will then need an enlarged operator basis, nonlocal channel mixing, or
a fully frequency-dependent formulation.  A negative result would still identify
why the proposed simplification fails.
